\section{Conclusion}
\noindent In this work, we have systematically analyzed the significant discrepancies between synthetic pretraining datasets and real-world pedestrian retrieval datasets, largely attributed to variations in illumination, pose, and other visual characteristics. 
To mitigate these domain shifts, we introduce a unified framework that tackles the domain gap at both the image and region levels. At the image level, we employ Domain-aware Diffusion (DaD) to reduce the stylistic and distributional divergence between the synthetic pretraining data and the real-world target dataset. This approach facilitates the effective utilization of synthetic data for pretraining, thereby enhancing the generalization capabilities of pedestrian retrieval models to real-world conditions.
At the region level, we augment the dataset by integrating region-phrase annotations with image-text pairs, thereby strengthening the semantic correspondence between visual content and descriptive text. We propose a Multi-granularity Relation Alignment (MRA) approach to fully leverage these enriched image-text pairs. MRA ensures comprehensive alignment between pedestrian images and corresponding text across different granularities, from coarse to fine, thus addressing the fine-grained discrepancies often challenging cross-domain retrieval.
The effectiveness and robustness of our dual-level adaptation approach, encompassing DaD and MRA, have been thoroughly validated through extensive experiments on three benchmark pedestrian retrieval datasets, \ie, CUHK-PEDES, ICFG-PEDES, and RSTPReid. 
Despite the promising results, we acknowledge the limitations of our approach. First, although DaD generates a large volume of domain-aligned images, it still inherits the inherent drawbacks of diffusion models, such as oversaturation in some generated images. Second, the captions for these images are produced by the BLIP2 model, a process that may introduce and propagate the inherent biases present in BLIP2.
Looking forward, our future work will explore several promising directions. A primary focus will be to extend this research from images to video-based person retrieval, investigating feasible solutions in more complex and realistic scenarios. This includes studying the impact of various human behaviors, especially abnormal activities~\cite{yang2025beyond}, on retrieval performance. We believe tackling these dynamic and complex settings is a crucial step toward real-world application.


\noindent\textbf{Acknowledgments.}
This work is supported by National Key Research and Development Program of China (2023YFC3321600).
Z. Zheng acknowledges supports from Guangdong Basic and Applied Basic Research Foundation 2025A1515012281, the University of Macau MYRG-GRG2024-00077-FST-UMDF, the University of Macau Advanced Research Institute in Hengqin, and the Macao Science and Technology Development Fund Grant FDCT/0043/2025/RIA1.

